%%%%%%%%%%%%%%%%%%%%%%%%%%%%%%%%%%%%%%%%%
% Seismica Submission Template
% Falcon 9 fireball paper working outline -- version 8
%%%%%%%%%%%%%%%%%%%%%%%%%%%%%%%%%%%%%%%%%

\documentclass[]{seismica}
\usepackage{tabularx}

% All manuscript graphics are stored in the figures subdirectory.
\graphicspath{{figures/}}

\title{Near-Field Seismo-Acoustic Observations of the 2016 Falcon 9 Explosion at Kennedy Space Center}
\shorttitle{Near-field observations of the 2016 Falcon 9 explosion}

% TODO: add ORCID identifiers and confirm corresponding-author address.
\author[1]{Glenn Thompson\thanks{Corresponding author: EMAIL TO BE CONFIRMED}}
\author[2]{Robert G. Brown}
\author[1]{Alexandra Farrell}
\author[2]{John Kiriazes}
\author[1]{Stephen R. McNutt}
\author[1]{Christopher Mehta}
\affil[1]{School of Geosciences, University of South Florida, Tampa, Florida, USA}
\affil[2]{NASA Kennedy Space Center, Florida, USA (current affiliations to be confirmed)}

% TODO: complete and verify CRediT contributions.
\credit{Conceptualization}{Glenn Thompson, Stephen R. McNutt}
\credit{Methodology}{Glenn Thompson}
\credit{Software}{Glenn Thompson}
\credit{Formal Analysis}{Glenn Thompson}
\credit{Investigation}{Glenn Thompson, Robert G. Brown, Alexandra Farrell, John Kiriazes, Stephen R. McNutt, Christopher Mehta}
\credit{Visualization}{Glenn Thompson}
\credit{Writing - Original draft}{Glenn Thompson}
\credit{Writing - Review \& Editing}{All authors}

\begin{document}

\makeseistitle{
  \begin{summary}{Abstract}
  On 1 September 2016, a Falcon~9 launch vehicle was destroyed during
  preparations for a pre-launch static-fire test at Space Launch Complex~40
  (SLC-40), Cape Canaveral Air Force Station (CCAFS; now Cape Canaveral Space
  Force Station), Florida. The accident was recorded approximately
  \SI{1.4}{\kilo\metre} away by a colocated broadband seismometer and
  three-element infrasound array operated by the University of South Florida
  near the Astronaut Beach House on adjacent Kennedy Space Center (KSC)
  property, providing an unusually close calibrated seismo-acoustic record of
  a major technological explosion.
  %
  The observations were analyzed as an independent forensic record of the
  accident: to test for detectable energetic precursors, reconstruct the
  chronology, quantify event amplitudes, determine whether signals originated
  consistently from the launch complex, and evaluate their propagation
  characteristics. Manual association identified 154 candidate infrasound events spanning
  approximately \SI{26.64}{\minute}; 153 yielded a valid canonical
  positive-compression/negative-rarefaction measurement, while the principal
  explosion was retained as a complex noncanonical event. No
  anomalous seismic or infrasonic precursor was detected during the preceding
  \SI{37}{\hour}. The initial upper-stage signal had an aligned-stack positive pressure of
  \SI{23}{\pascal} and peak-to-peak pressure of \SI{44}{\pascal};
  approximately \SI{3.67}{\second} later, the principal explosion generated
  an aligned-stack positive pressure of \SI{1.353}{\kilo\pascal}, a
  peak-to-peak pressure of \SI{1.437}{\kilo\pascal}, and the strongest
  ground motion of the sequence. A prolonged interval of
  continuous signal followed the initial phase, but its physical origin
  remains uncertain.
  %
  Array solutions consistently pointed toward SLC-40 and exhibited stable
  apparent propagation speeds despite more than two orders of magnitude
  variation in pressure amplitude. Seismic records were dominated by
  ground-coupled airwaves, and classical N-wave morphology was not observed.
  A single near-field seismo-acoustic array therefore provided an independent
  reconstruction of a complex launch-vehicle failure without reliance on
  proprietary engineering telemetry.
  \end{summary}
  %
  \begin{summary}{Non-technical summary}
  A seismic and low-frequency sound monitoring station approximately
  \SI{1.4}{\kilo\metre} from the launch pad recorded the 2016 Falcon~9
  accident at SLC-40. The launch pad was at Cape Canaveral Air Force Station,
  while the monitoring station was located on adjacent Kennedy Space Center
  property. The recordings provide an independent account of what happened:
  whether any detectable energetic signal preceded the visible failure, how
  many subsequent events occurred, how large they were, where they originated,
  and how quickly their pressure waves travelled. No unusual precursor was
  detected during the preceding \SI{37}{\hour}. The accident was not a single explosion, but a sequence of 154 associated
  candidate pressure events extending over approximately
  \SI{26.64}{\minute}; 153 were represented by the standard short-pulse
  measurement algorithm. By comparing arrival times and waveforms at
  three infrasound sensors, we show that the events consistently originated
  from the direction of the launch complex. The study demonstrates how a
  small, continuously recording geophysical station can provide an independent
  forensic record of a complex launch-vehicle failure.
  \end{summary}
}


\section{Introduction}
\label{sec:introduction}

Rocket launches generate some of the strongest routinely occurring
anthropogenic seismic and infrasonic signals on Earth. During nominal launches,
intense acoustic radiation is produced principally by the interaction of rocket
exhaust with the atmosphere and launch infrastructure, and these airwaves can
couple efficiently into the shallow ground. Rocket-generated signals have been
studied for both operational monitoring and geophysical research, including
global surveys of infrasonic signatures from routine space launches
\citep{pilger_rocket_2021}. Catastrophic launch failures remain rare,
particularly when recorded at close range by calibrated seismic and infrasound
instruments.

On 1 September 2016, a Falcon~9 launch vehicle was destroyed during
preparations for a pre-launch static-fire test at Space Launch Complex~40
(SLC-40), Cape Canaveral Air Force Station (CCAFS; now Cape Canaveral Space
Force Station), Florida. The vehicle and its AMOS-6 payload were lost, and the
launch complex sustained extensive damage. The accident was documented by
several publicly available videos and was subsequently investigated by SpaceX,
providing a high-energy technological explosion with known source location,
timing, and geometry for which the geophysical record can be compared with
independent visual observations and a subsequently released engineering
chronology.

In the immediate aftermath, the USF recordings were examined as an independent
forensic record of the accident. The principal questions were whether any
detectable energetic precursor preceded the visible failure; how many discrete
events occurred and when; how large the principal events were; whether the
signals originated consistently from SLC-40; and whether their apparent
propagation speeds were compatible with atmospheric airwaves under the observed
meteorological conditions. These questions motivated the original analysis
conducted for SpaceX and the subsequent scientific investigation.

At the time of the accident, the University of South Florida (USF) was
operating the BCHH seismo-acoustic array near the Astronaut Beach House on
Kennedy Space Center property, approximately 1.4~km north-northeast of SLC-40.
The installation comprised a colocated broadband seismometer and three-element
infrasound array and formed part of a broader USF program that used rocket
launches to evaluate seismo-acoustic instrumentation between volcanic field
deployments, train students in field operations, and investigate near-field
rocket-generated wavefields. The BCHH array fortuitously recorded the complete
accident sequence with high signal-to-noise ratio, without reliance on vehicle
telemetry or a triggered recording system.

Previous seismo-acoustic studies of large technological explosions have
generally relied on regional observations. Examples include the 2005 Buncefield
oil-depot explosion \citep{ottemoller_2008,ceranna_2009,green_2008}, the 2014
Moerdijk chemical-plant explosions \citep{smink_2019}, and the 2020 Beirut
ammonium-nitrate explosion \citep{pilger_beirut_2021,kim_pasyanos_2022}.
Atmospheric propagation modelling was central to interpreting arrival times,
amplitudes, source location, and explosive yield in those studies. By contrast,
the 1.4-km source distance places the BCHH array in the near field, allowing
direct-arrival analysis with comparatively little influence from atmospheric
stratification or ducting, while still requiring explicit treatment of local
wind and finite-range geometry. This proximity allows the chronology and
relative strength of more than 150 individual impulsive events to be resolved
and compared with independent video observations. The dataset therefore
provides an unusually detailed view of the evolution of a catastrophic
launch-vehicle failure.

The objectives of this study are to:
\begin{enumerate}
  \item test the continuous seismic and infrasound record for detectable
        energetic precursors to the visible failure;
  \item construct a quantitative chronology of the accident and determine the
        number and temporal distribution of discrete impulsive events;
  \item measure the acoustic and seismic amplitudes of the principal events
        and characterize their associated wavefields;
  \item determine whether the impulsive arrivals originated consistently from
        the direction of SLC-40 and whether their apparent propagation speeds
        were compatible with the surveyed geometry and observed atmospheric
        conditions; and
  \item evaluate the value and limitations of near-field seismo-acoustic
        observations for independent forensic analysis of launch-vehicle
        accidents and other complex technological explosions.
\end{enumerate}

This study does not attempt to determine the internal engineering mechanism
that initiated the vehicle failure. In particular, the absence of a detectable
energetic precursor cannot exclude non-energetic or seismo-acoustically
undetectable processes. Instead, the study focuses on what can be inferred
directly from the near-field geophysical observations and on how those
observations complement the publicly available visual record and the
subsequently released failure investigation.


\section{Instrumentation and data}
\label{sec:data}

\subsection{BCHH seismo-acoustic array}

The BCHH array was initially deployed on 24 February 2016 as part of a
three-station USF network on Kennedy Space Center property, but two stations
were removed in May 2016. The network was established to evaluate and
recalibrate seismic and infrasound instruments returned from the 2015
Sakurajima experiment, during which anomalous amplitude responses had been
identified. It subsequently developed into a longer-term program in which
equipment was kept in continuous operational use between volcanic field
campaigns. This supported routine instrument-performance checks, student
training in equipment deployment and maintenance, teaching, and research on
rocket-generated seismic and acoustic wavefields. The station and channel history, 
including sensor positions and response metadata, is preserved in the network 
metadata and StationXML files used during waveform processing.

The 1 September 2016 Falcon~9 accident at SLC-40, Cape Canaveral Air Force
Station (CCAFS), was recorded at BCHH near the Astronaut Beach House on
adjacent Kennedy Space Center (KSC) property
(Figure~\ref{fig:site_array}). Located approximately 1.42~km
north-northeast of SLC-40, BCHH provided favorable geometry for recording
both routine launches and, fortuitously, the accident. At the time of the
event, the array comprised a three-component Nanometrics Trillium Compact
Posthole broadband seismometer near the center of a three-element infrasound
array with an aperture of approximately 30~m. The seismometer was buried
approximately 0.5~m below ground level to improve thermal stability and reduce
wind-induced noise.

The three infraBSU differential-pressure sensors were mounted individually on
short stakes approximately 20--30~cm above ground level in an approximately
equilateral triangular configuration around the seismometer. All six channels
were digitized continuously at 250 samples~s$^{-1}$ using a GPS-synchronized
Nanometrics Centaur recorder. 

\begin{figure*}[t]
  \centering
  \includegraphics[width=0.96\textwidth]{fig01_slc40_bchh_location.pdf}
  \caption{Location and geometry of the BCHH seismo-acoustic array.
  (a) Map of the southern Kennedy Space Center--Cape Canaveral launch area
  showing the USF BCHH array on Kennedy Space Center property, approximately
  1.4~km north-northeast of Space Launch Complex~40 (SLC-40) at Cape
  Canaveral Air Force Station; Space Launch Complex~41 is shown for
  geographic reference.
  (b) Geometry of the BCHH array, comprising a three-component broadband
  seismometer near the center of a three-element triangular infrasound array
  with an aperture of approximately 30~m. Sensor locations were determined by
  differential-GPS survey.}
  \label{fig:site_array}
\end{figure*}


\subsection{Surveyed geometry}

The sensor positions were determined by differential-GPS survey on 2 November
2016, yielding estimated horizontal uncertainties of approximately 2--3~cm
(Table~S2). Using the surveyed receiver coordinates and the known launch-pad
location of the Falcon~9 vehicle, horizontal source-to-sensor distances for
the three infrasound sensors ranged from 1406.3 to 1438.0~m, while their
geometric back azimuths ranged from 198.73$^\circ$ to 199.91$^\circ$
(Table~S3). These coordinates formed the basis for
the differential-travel-time corrections and planar- and finite-range array
analyses described in Section~\ref{sec:methods}.


\subsection{Instrument responses and pressure calibration}

The broadband-seismometer response was removed using the manufacturer
calibration for the Trillium Compact Posthole--Centaur instrument chain. The
infraBSU channels were converted to pressure using empirical, sensor-specific
field calibration factors derived from comparative tests across the set of 25
infraBSU sensors operated by USF. Available calibration records indicate that
these experiments included colocated lift tests and impulsive tests such as
door slams, allowing relative differences in sensor sensitivity to be
quantified. The adopted BCHH calibration factors are summarized in
Table~\ref{tab:instrument_calibration}; their implementation during response
correction is described in Section~\ref{sec:methods}.

The original analysis assigned an approximate absolute pressure uncertainty of
$\pm20\%$. Because the basis for this value is not encoded in the modern
processing notebooks, it will be verified against the original calibration
records before final submission. Timing uncertainty reflects the 4-ms sampling
interval, GPS synchronization, and manual arrival picking. For clear impulsive
arrivals, absolute pick uncertainties are estimated to be better than
approximately 0.02~s.

\begin{table}
\centering
\caption{Instrument channels, sampling rate, and calibration factors used for
response correction. The infraBSU values are empirical field sensitivities.}
\label{tab:instrument_calibration}
\begin{tabular}{lll}
\hline
Instrument/channel & Sampling rate & Sensitivity \\
\hline
Trillium Compact HHZ/HHN/HHE
  & 250~Hz
  & $3.0172\times10^8$ counts $(\mathrm{m\,s^{-1}})^{-1}$ \\
infraBSU HD1
  & 250~Hz
  & 8.10 counts~Pa$^{-1}$ \\
infraBSU HD2
  & 250~Hz
  & 0.690 counts~Pa$^{-1}$ \\
infraBSU HD3
  & 250~Hz
  & 10.0 counts~Pa$^{-1}$ \\
\hline
\end{tabular}
\end{table}


\subsection{Meteorological observations}

Meteorological observations were obtained from the KSC weather-tower network
during the accident sequence. The network comprised 68 instrumented towers
spanning measurement heights of approximately 4--150~m above ground level.
Available variables included wind speed, wind direction, air temperature,
relative humidity, and barometric pressure at 5-min intervals between 13:05
and 13:35~UTC. These observations provided unusually dense spatial and vertical
sampling of atmospheric conditions during the accident and were used to
estimate the still-air sound speed, the along-path wind correction, and the
expected acoustic travel time between SLC-40 and BCHH
(Section~\ref{sec:methods}).


\subsection{Publicly available video and audio}

The seismic and infrasound observations were compared with publicly available
video recordings of the accident. These recordings provide an independent
visual chronology and allow selected geophysical observations to be related to
visible stages including the initial flash, growth of the fireball, structural
collapse, descent of vehicle or payload components, and subsequent energetic
activity. Public availability makes the source material independently
inspectable, although several recordings lack reliable absolute timestamps and
their camera positions, frame rates, and timing uncertainties require further
documentation.

The videos were used principally to interpret and label the geophysical
sequence. Event association, waveform amplitudes, array measurements, and the
other quantitative catalogue products were derived from the BCHH seismic and
infrasound recordings. Proprietary SpaceX footage may have been viewed during
the original forensic work, but no quantitative result reported here will rely
on material that cannot be independently inspected or redistributed; any
historical influence of such footage on an interpretation will be disclosed
explicitly. The procedures and limitations of video and audio
synchronization are described in Section~\ref{sec:methods}.


\section{Methods}
\label{sec:methods}

\subsection{Analysis overview TO DO}

The analysis was organized around the principal forensic questions posed
following the accident: whether any detectable energetic precursor preceded
the visible failure; how many discrete events occurred and when; how their
acoustic and seismic amplitudes varied; whether they originated consistently
from SLC-40; and whether their apparent propagation speeds were compatible with
direct acoustic propagation under the observed atmospheric conditions.

The technical workflow comprised:
(1) correction of instrument responses and conversion to physical units;
(2) compilation of surveyed source--receiver geometry and meteorological
constraints;
(3) manual identification and reduced-time association of candidate airwave
arrivals;
(4) measurement of pressure-pulse morphology, amplitude, and signal-to-noise
ratio;
(5) estimation of back azimuth and apparent propagation speed using
three-channel planar-wave coherence analysis;
(6) assessment of finite-range effects using a horizontal circular-wavefront
model;
(7) measurement and comparison of acoustic pressure and seismic peak ground
velocity;
(8) calculation of time-dependent three-component polarization attributes and
particle-motion measures;
(9) estimation of acoustic and seismic energies and equivalent magnitudes
using the original MATLAB workflow;
(10) comparison of the geophysical chronology with publicly available video
and audio recordings; and
(11) generation of a synchronized forensic visualization combining video,
waveforms, and polarization attributes.

Modern waveform processing, event association, amplitude measurement, and
array analysis were implemented in Python and are described directly from the
processing code. Python-derived results take precedence wherever an equivalent
modern product exists. Energy and equivalent-magnitude calculations that have
not yet been reimplemented in Python are retained provisionally from the
original MATLAB workflow and will be replaced or independently verified when a
Python implementation is available. The legacy MATLAB/GISMO workflow also
provided the original synchronized video, waveform, and polarization display,
but interpretations derived from that product remain provisional until
reproduced using the documented Python workflow. Continuous waveform data and
metadata were originally archived in a CSS3.0 database using Antelope
\citep{brtt_antelope}, manual arrival picks were made using
\texttt{dbpick}, and the original MATLAB analysis used GISMO
\citep{thompson_gismo_2017}. The canonical Python analysis interval extended from 180~s
before to 1800~s after a reference source time of 13:07:12.080~UTC assigned to
the initial upper-stage event. All six waveform channels were sampled
continuously at 250~Hz. The polarization analysis and synchronized visualization are being
reimplemented in Python and will become authoritative once their outputs have
been validated against the archived MATLAB products.


\subsection{Instrument calibration and response correction}

The three seismic components were corrected using the manufacturer response
for the Nanometrics Trillium Compact Posthole seismometer and Centaur digitizer
combination. The combined sensitivity was
$3.0172\times10^8$ counts $(\mathrm{m\,s^{-1}})^{-1}$. The Centaur was
operated with a 40-V peak-to-peak input range, corresponding to a digitizer
sensitivity of $4.0\times10^5$ counts~V$^{-1}$ and an implied seismometer
sensitivity of approximately 754~V~$(\mathrm{m\,s^{-1}})^{-1}$.

The three infraBSU differential-pressure sensors were converted independently
from digital counts to pressure using empirical, sensor-specific field
calibration factors of 8.10, 0.690, and 10.0 counts~Pa$^{-1}$ for HD1, HD2,
and HD3, respectively. These factors were derived from comparative experiments
in which colocated sensors were exposed to common pressure changes and
impulsive sources. They should therefore be regarded as empirical field
calibrations rather than laboratory-traceable absolute calibrations. The
precise combination of tests used to derive the adopted factors, and the basis
for the previously assigned absolute pressure uncertainty of approximately
$\pm20\%$, will be verified from the original calibration records.

For each infrasound channel, the median count level between 20 and 2~s before
the initial explosion was removed before response correction. The empirical
pressure sensitivities were incorporated into complete instrument-response
representations while retaining the nominal poles, zeros, and digitizer
stages. Seismic traces were demeaned and linearly detrended. All traces received
a 5\% cosine taper before frequency-domain response removal.

A four-corner prefilter of 0.02, 0.05, 80, and 100~Hz and a water level of
60~dB were used to stabilize the response inversion. This prefilter was applied
only during response removal. No additional Butterworth high-pass filter was
applied during event association, amplitude measurement, or array analysis.
Corrected seismic traces were expressed as ground velocity in
m~s$^{-1}$, and corrected infrasound traces as differential pressure in Pa.


\subsection{Source--receiver geometry}

Sensor coordinates were obtained from the differential-GPS survey of BCHH and
transformed to UTM Zone~17N for array calculations. Horizontal
source--receiver distances and geographic azimuths were also calculated
independently on the WGS84 ellipsoid. The central broadband seismometer, used
as the array reference position, was located at approximately 541820.4~m~E and
3160866.5~m~N. The reference horizontal distance from SLC-40 to the central broadband
seismometer was 1417.9~m, and the geometric back azimuth from that nominal
array-reference position toward SLC-40 was 199.41$^\circ$, measured clockwise
from north.

The three infrasound sensors formed an array with an aperture of approximately
30~m. Their horizontal distances to SLC-40 were 1438.0~m (HD1), 1406.3~m
(HD2), and 1410.9~m (HD3), with corresponding geometric back azimuths of
198.95$^\circ$, 199.91$^\circ$, and 198.73$^\circ$, respectively. Individual
source-to-sensor distances were retained when calculating expected
differential acoustic travel times rather than treating the array as a single
receiver. HD2 was used as the infrasound reference channel for reduced-time
association and waveform alignment; consequently, its geometric direction to
SLC-40 differs slightly from the nominal array-center direction.

Approximate source-time labels for selected stages of the accident were
obtained from the video chronology. These labels assisted interpretation of the
initial upper-stage event, the principal explosion, and later
capsule- or payload-related activity, but were not required to construct the
general event catalogue, which was based on observed infrasound arrivals. A
possible capsule impact has not yet been confirmed directly in the raw seismic
waveforms and will be reassessed using polarization attributes regenerated in
the Python workflow. Formal uncertainties for the video-derived source-time
labels have not yet been established, and all but the best-constrained onset
time should therefore be regarded as approximate.

\subsection{Legacy chronology values and reconciliation}
\label{sec:methods_legacy_chronology}

The archived MATLAB analysis contains several chronology values that are useful
for comparison but are not adopted as final measurements until they have been
reproduced or reconciled with the Python workflow. These include a weak seismic
feature at approximately 13:07:13.65~UTC, a first infrasound pulse approximately
2.28~s later, a largest-event arrival near 13:07:19.43~UTC, payload- or
capsule-related activity near 13:07:24.5~UTC, a possible seismic onset near
13:07:25.6~UTC, and the beginning of the prolonged continuous signal near
13:08:40~UTC. The legacy high-quality array subset also yielded a median
apparent propagation speed of approximately 358.6~m~s$^{-1}$.

These values mix observed arrival times, inferred source times, video-derived
labels, and legacy processing conventions. They will therefore be reconciled
against the Python reduced-time catalogue, the canonical reference time of
13:07:12.080~UTC, the provisional principal-explosion source time, and the
final video synchronization before they are used in tables, figures, or
quantitative interpretation. Where modern Python and legacy MATLAB values
differ, the Python result will be reported as primary and the legacy value will
be retained only for documented historical comparison.


\subsection{Meteorological correction}

Meteorological observations from the Kennedy Space Center tower network between
13:05 and 13:35~UTC were used to estimate the effective acoustic propagation
speed between SLC-40 and BCHH. The selected dataset contained 84 temperature
and 84 relative-humidity observations, with mean conditions of approximately
26.36$^\circ$C and 83.87\% relative humidity.

The still-air sound speed was calculated as
\begin{equation}
  c_0 = 331.3 + 0.606T + 1.26h,
  \label{eq:sound_speed}
\end{equation}
where $T$ is temperature in degrees Celsius and $h$ is relative humidity
expressed as a fraction. This yielded
$c_0=348.33~\mathrm{m\,s^{-1}}$.

Wind observations were resolved into east and north components, averaged
within 10-m height intervals, and linearly interpolated through the lowest
65~m of the atmosphere. The height-averaged wind component parallel to the
SLC-40--BCHH propagation path was added to the thermodynamic sound speed:
\begin{equation}
  c_{\mathrm{eff}} =
  c_0 + \mathbf{u}\cdot\hat{\mathbf{r}}.
  \label{eq:effective_speed}
\end{equation}

The mean along-path wind component was $+2.65~\mathrm{m\,s^{-1}}$, giving
$c_{\mathrm{eff}}=350.98~\mathrm{m\,s^{-1}}$. The corresponding travel time
over 1417.9~m was 4.040~s, approximately 0.031~s shorter than the still-air
estimate. Because the propagation distance was only about 1.4~km, direct
atmospheric arrivals dominated and long-range ducting effects were not
considered. The meteorologically estimated value of
350.98~m~s$^{-1}$ is distinct from the approximately
366.5--366.8~m~s$^{-1}$ empirical speed used only to align one supplementary
video/audio display.

\subsection{Manual arrival picking and baseline correction}

Continuous waveforms were visually reviewed using Antelope \texttt{dbpick}. Candidate events were initially identified as impulsive pressure arrivals on the infrasound channels. Only non-deleted phase-$N$ picks on HD1, HD2, and HD3 were retained, and the original pick precision was preserved when exporting the arrivals to the Python workflow.

Manual review was adopted because the accident sequence occupied approximately 26~min within a modest six-channel dataset and contained closely spaced, overlapping arrivals spanning a wide range of amplitudes and waveform shapes. For this dataset, manual inspection provided a consistent reference catalogue without the iterative tuning and validation required for a new automated detector. Local seismic transients without corresponding infrasound arrivals were excluded from the airwave catalogue.

To address the forensic question of whether any detectable energetic precursor
preceded the visible failure, the complete available record from 31 August
through 1 September 2016 was also examined outside the main accident sequence.
This interval included approximately 37~h before the first accident-related
signal and 10.5~h after the final catalogued event. The search was intended to
identify anomalous impulsive seismic or infrasonic signals and cannot exclude
non-energetic, aseismic, or otherwise undetectable initiating processes.

Slow baseline variations on the response-corrected infrasound channels were removed using a centered, noncausal moving median with a 251-sample window, corresponding to approximately 1~s at 250~Hz. The estimated median baseline was subtracted from each pressure trace. Seismic traces were not subjected to this baseline operation.

The moving-median procedure was selected after comparison with unfiltered waveforms, conventional high-pass-filtered traces, local linear baselines, and moving-median windows of several durations. High-pass filtering enhanced some weak impulsive arrivals but distorted broader pressure pulses, whereas local linear baselines were occasionally controlled by a small number of samples and produced implausibly steep trends. The 1-s moving median provided the most stable baseline estimate while preserving the short pressure impulses. Although a centered median is nonlinear and therefore not formally a zero-phase filter, it does not impose the systematic delay of a one-sided causal smoother.

\subsection{Reduced-time event association}

Before event association, each accepted pick was corrected for its expected differential travel time relative to HD2:
\begin{equation}
  t_{i,\mathrm{red}} = t_i - \frac{r_i-r_{\mathrm{HD2}}}{c_{\mathrm{eff}}},
  \label{eq:reduced_time}
\end{equation}
where $r_i$ is the SLC-40--sensor distance and $c_{\mathrm{eff}}=350.98~\mathrm{m\,s^{-1}}$. The resulting corrections were approximately 0.090~s for HD1, 0~s for HD2, and 0.013~s for HD3. These corrections compensated for the known source--receiver geometry only during preliminary association; final propagation direction and apparent speed were estimated independently from the waveforms.

Reduced picks were sorted chronologically. Beginning with the earliest unassigned pick, all picks falling within the following 0.040~s were treated as a candidate cluster. At most one pick per channel was retained, with the pick closest to the cluster seed selected when multiple picks occurred on the same channel. A cluster was accepted as a candidate event if it contained picks from at least two distinct infrasound channels. Picks that failed this criterion remained unassociated rather than being forced into an event.

For each accepted candidate, a six-channel waveform segment was extracted from 0.20~s before to 0.30~s after the reduced event time. The provisional travel-time corrections were not written into these waveform streams; the original observed sample times were retained. Application of the association procedure produced the candidate catalogue used in the subsequent waveform and array analyses.

\subsection{Acoustic pulse measurements and event-quality tiers}

For initial pulse characterization, the three infrasound traces were shifted to the HD2 reference time using the meteorologically predicted differential delays and combined using a sample-by-sample median. Within a search interval from $-0.04$ to $+0.25$~s relative to the reduced event time, the algorithm sought a positive compression followed by a negative pressure excursion. The negative peak was required to occur 0.010--0.180~s after the positive peak. When several peak pairs satisfied these criteria, the pair with the largest peak-to-peak pressure was retained.

For positive and negative peaks at times $t_+$ and $t_-$, provisional one-cycle bounds were defined as
\begin{equation}
  t_{\mathrm{start}} = t_+ - 0.5(t_- - t_+)
  \label{eq:pulse_start}
\end{equation}
and
\begin{equation}
  t_{\mathrm{end}} = t_- + 0.5(t_- - t_+).
  \label{eq:pulse_end}
\end{equation}
Although one configuration option proposed refining these boundaries to the return-to-noise points, that refinement was not implemented in the adopted measurement function and is not assumed here.

Noise was measured from $-0.40$ to $-0.10$~s relative to the event using the robust scale
\begin{equation}
  \sigma_{\mathrm{robust}} = 1.4826\,\mathrm{median}\left(\left|p-\mathrm{median}(p)\right|\right).
  \label{eq:robust_noise}
\end{equation}
The provisional stack signal-to-noise ratio was the maximum absolute pressure within the event bounds divided by this noise scale.

The positive--negative pulse morphology was used as a measurement-quality criterion rather than as the defining requirement for a physical event. In particular, the principal explosion was retained despite failing the simple one-cycle morphology test because it was an exceptionally strong, coherent, multichannel signal with a more complex pressure history.

Candidate labels based on channel support, reduced-pick span, pressure amplitude, and signal-to-noise ratio were retained as descriptive quality indicators. They were not interpreted as distinct physical source classes.

\subsection{Eligibility for detailed array analysis}

Two operational threshold sets were used to identify events suitable for array analysis. The permissive analysis set required a median-stack peak-to-peak pressure of at least 10~Pa, a stack signal-to-noise ratio of at least 1.5, signal support on at least two infrasound channels, and a valid positive--negative pulse measurement. The principal explosion was explicitly retained despite failing the short-pulse morphology requirement.

A more conservative core subset required a median-stack peak-to-peak pressure of at least 25~Pa and a stack signal-to-noise ratio of at least 3, together with the same channel-support requirement. These thresholds were configurable operational divisions along continuous amplitude and quality distributions; they were not interpreted as natural boundaries between event populations. All filtering decisions, including reasons for rejecting candidates, were retained in the analysis products.

\subsection{Planar-wave array analysis}

Back azimuth and apparent horizontal propagation speed were estimated jointly by shifting all three infrasound waveforms according to trial plane-wave models. Back azimuth $\theta$ was defined as the direction from the array toward the source, measured clockwise from north. For a trial back azimuth and apparent speed $v$, the horizontal propagation vector toward the array was
\begin{equation}
  \mathbf{k} = [-\sin\theta,-\cos\theta],
  \label{eq:propagation_vector}
\end{equation}
and the predicted relative delay at sensor $i$, relative to HD2, was
\begin{equation}
  \Delta t_i = \frac{(\mathbf{x}_i-\mathbf{x}_{\mathrm{HD2}})\cdot\mathbf{k}}{v}.
  \label{eq:plane_delay}
\end{equation}

For each trial model, the three waveforms were shifted by interpolation and trimmed to their common temporal overlap. Every point on the search surface was therefore evaluated using the same waveform duration, avoiding the unequal-overlap bias that can favor zero lag in conventional cross-correlation calculations.

Each aligned trace was median-centered and normalized by its root-mean-square amplitude. Coherence was evaluated using both three-channel semblance,
\begin{equation}
  S = \frac{\sum_t\left[\sum_i x_i(t)\right]^2}{N\sum_t\sum_i x_i^2(t)},
  \label{eq:semblance}
\end{equation}
and the mean of the three pairwise correlation coefficients. The mean correlation $\overline{\rho}$ was mapped from $[-1,1]$ to $[0,1]$, and the combined score was
\begin{equation}
  C = 0.60S + 0.40\left(\frac{\overline{\rho}+1}{2}\right).
  \label{eq:coherence_score}
\end{equation}

The coarse search covered back azimuths from 180$^\circ$ to 220$^\circ$ in 0.5$^\circ$ increments and apparent speeds from 280 to 440~m~s$^{-1}$ in 2~m~s$^{-1}$ increments. Each coarse maximum was refined over $\pm2^\circ$ in 0.1$^\circ$ increments and $\pm20$~m~s$^{-1}$ in 0.5~m~s$^{-1}$ increments. Most events were scored over the interval $-0.04$ to $+0.18$~s relative to the reduced event time; the principal explosion used an extended interval ending at $+0.35$~s.

Pairwise correlations and residual time lags were calculated after application of the best-fitting model. The lag-closure quantity
\begin{equation}
  \tau_{12}+\tau_{23}+\tau_{31}
  \label{eq:lag_closure}
\end{equation}
was retained as a consistency diagnostic. The relative performance of HD1 was evaluated by comparing the HD2--HD3 correlation with the mean correlation of the two pairs containing HD1.

Solution stability was tested by shifting the scoring window by $-0.008$, 0, and $+0.008$~s and scaling its duration by factors of 0.9, 1.0, and 1.1. Standard deviations of the resulting back azimuths and speeds were retained as empirical stability measures. Widths of the coherence surface at score reductions of 0.01 and 0.02 were also calculated, but these widths were treated as descriptive measures rather than formal confidence intervals.

Composite coherence surfaces were formed from the individual coarse-search surfaces for both the full analysis set and the core subset. Arithmetic-mean, median, maximum-score-weighted, and individually normalized mean composites were retained to evaluate the population-level consistency of the solutions.

\subsection{Finite-range circular-wavefront analysis}

The adequacy of the plane-wave approximation was assessed using exact horizontal source--receiver distances. For a trial source position $\mathbf{x}_s$ and propagation speed $v$, the predicted delay at sensor $i$, relative to HD2, was
\begin{equation}
  \Delta t_i = \frac{|\mathbf{x}_i-\mathbf{x}_s|-|\mathbf{x}_{\mathrm{HD2}}-\mathbf{x}_s|}{v}.
  \label{eq:circular_delay}
\end{equation}
Because receiver elevations were not included in the adopted geometry, this was a two-dimensional circular-wavefront model rather than a full three-dimensional spherical-wave model.

The primary finite-range calculation fixed the source at SLC-40 and searched propagation speeds from 280 to 440~m~s$^{-1}$ in 0.5~m~s$^{-1}$ increments. Coherence was evaluated using the same waveform windows, normalization, and scoring metric as the planar-wave analysis.

Sensitivity to source position was examined in coordinates oriented parallel and perpendicular to the SLC-40--BCHH path. Along-range offsets from $-500$ to $+500$~m and cross-range offsets from $-150$ to $+150$~m were evaluated while profiling over the complementary coordinate. These calculations were interpreted as resolution tests rather than independent source locations. With three sensors, only two independent relative delays are available, so source easting, source northing, and propagation speed cannot all be determined independently. Open along-range coherence ridges were therefore interpreted as a fundamental range-resolution limitation of the array geometry.

\subsection{Acoustic and seismic amplitude measurements}

Acoustic and seismic amplitudes were measured for each candidate event using the saved baseline-corrected waveform segments. Event intervals were taken from the positive--negative pulse measurements, except for the principal explosion, for which a fixed interval from $-0.04$ to $+0.35$~s was used.

For each seismic component, the median ground velocity from $-0.18$ to $-0.05$~s was removed. Component peak ground velocity was the largest absolute velocity within the event interval. Three-component vector peak ground velocity was calculated as
\begin{equation}
  \mathrm{PGV}_{\mathrm{vec}} = \max_t\sqrt{v_E^2(t)+v_N^2(t)+v_Z^2(t)}.
  \label{eq:vector_pgv}
\end{equation}
The maximum vector amplitude over the complete saved waveform segment was also retained to determine whether the acoustically defined event interval captured the largest seismic motion.

For catalogue measurements, the three infrasound traces were aligned to HD2 using the meteorologically predicted differential delays and combined by a sample-by-sample median. Positive pressure, negative pressure, and peak-to-peak pressure were measured within the same event interval. Acoustic signal-to-noise ratio was defined as median-stack peak-to-peak pressure divided by the robust pre-event noise scale.

The median-stack measurements provide a consistent catalogue metric but differ from measurements made separately on the three physical sensors. For the principal explosion, absolute positive overpressure was therefore characterized from the individual calibrated sensors, whereas the aligned median stack was used for comparison with the remaining catalogue. The earlier sensor-specific analysis gave positive peaks of approximately 1349--1557~Pa, with a median of 1435~Pa, and negative excursions of approximately $-179$ to $-223$~Pa. The present median-stack calculation gives approximately $+1353$~Pa, $-84$~Pa, and 1437~Pa peak-to-peak. These quantities represent different measurement procedures and should not be used interchangeably.

Where pressure was reduced to a reference distance of 1~km, spherical geometric spreading was assumed:
\begin{equation}
  p_{\mathrm{red},1\,\mathrm{km}} = p_{\mathrm{obs}}\frac{r}{1000~\mathrm{m}}.
  \label{eq:reduced_pressure}
\end{equation}
For a positive pressure peak of 1435~Pa at the BCHH distance, this gives a reduced positive pressure of approximately 2.03~kPa at 1~km.

Peak pressure level was calculated as
\begin{equation}
  L_{\mathrm{peak}} = 20\log_{10}\left(\frac{p_+}{20~\mu\mathrm{Pa}}\right).
  \label{eq:peak_pressure_level}
\end{equation}
This is a peak pressure level, not an RMS or time-averaged sound-pressure level. The positive peak pressure of approximately 1435~Pa corresponds to a peak level of approximately 157~dB re 20~$\mu$Pa.

A dimensional acoustic-to-seismic amplitude ratio was calculated as peak-to-peak pressure divided by vector peak ground velocity, with units of Pa per m~s$^{-1}$. This quantity was treated as an empirical measure of pressure--ground-motion coupling and not as an acoustic-to-seismic energy ratio.

\subsection{Acoustic--seismic amplitude regression}

Regression analysis was restricted to events with vector-peak-ground-velocity signal-to-noise ratio of at least 5, acoustic peak-to-peak signal-to-noise ratio of at least 5, and an event-window peak ground velocity equal to at least 70\% of the maximum peak ground velocity measured in the complete saved segment.

A free power-law relationship,
\begin{equation}
  p_{\mathrm{p2p}} = A\,\mathrm{PGV}^{b},
  \label{eq:power_law}
\end{equation}
was fitted by ordinary least squares after transforming both variables to base-10 logarithms. A 95\% descriptive interval for the exponent $b$ was estimated using 5000 nonparametric bootstrap resamples of the event population.

A proportional model,
\begin{equation}
  p_{\mathrm{p2p}} = K\,\mathrm{PGV},
  \label{eq:proportional_model}
\end{equation}
was also evaluated as a constant-ratio alternative. The two models were compared using root-mean-square residuals in base-10 logarithmic space. These regressions describe the observed relationship only; they do not account for measurement uncertainty in both variables, covariance among sensor-calibration errors, or censoring near the acoustic or seismic noise floors.

\subsection{Acoustic and seismic energy estimates}

Acoustic and seismic energies were estimated in the original MATLAB analysis using the formulations of \citet{johnson_aster_2005}, developed for partitioning acoustic and seismic energy from impulsive volcanic sources. Equivalent energy magnitudes were subsequently calculated using the relationship of \citet{kanamori_1977}, and the ratio of acoustic to seismic energy was used to characterize the relative partitioning of radiated energy between the atmosphere and ground. These calculations produced, for the principal explosion, first-order estimates of approximately 1030~MJ acoustic energy and 7.8~MJ seismic energy, corresponding to equivalent acoustic and seismic magnitudes of approximately 2.88 and 1.46, respectively.

These calculations have not yet been regenerated by the canonical Python notebook sequence and therefore remain products of the archived MATLAB workflow. The final version of this subsection should be checked directly against the original MATLAB routines to document the exact integration intervals, waveform quantities, density and wave-speed assumptions, propagation geometry, radiation assumptions, and treatment of the three sensors and seismic components. Until those implementation details are recovered, the resulting energies should be regarded as first-order estimates, most reliable for relative comparison and energy-partition interpretation rather than as precise source-yield determinations.

\subsection{Video and audio comparison}

Publicly available video recordings were used to identify major visible stages
of the accident and to assign approximate source-time labels to the initial
failure, the principal explosion, the descent and apparent ground contact of
the payload or capsule, and later capsule- or payload-related energetic
activity. Identifiable visual changes included the first flash, expansion of
the principal fireball, structural collapse, descent of vehicle components,
apparent ground contact, and subsequent explosions.

The quantitative event catalogue, pressure measurements, and array results were
derived independently from the BCHH geophysical data. Video observations were
used principally to interpret the physical processes associated with key
signals and to organize the accident chronology. In particular, a direct
seismic signal associated with the apparent payload or capsule impact has not
yet been confirmed in the raw three-component waveforms. The earlier impact
interpretation was based partly on polarization attributes calculated in the
legacy MATLAB workflow and will be reassessed after rectilinearity, planarity,
ellipticity, and particle-motion measures are regenerated in Python. Because
acoustic travel times from SLC-40 to the different cameras varied and several
videos lacked reliable absolute timestamps, synchronization relied primarily
on visible frames rather than on recorded audio.

For a supplementary opening-sequence visualization, audio extracted from one publicly available video was demeaned and band-pass filtered from 1 to 2~kHz using a four-pole zero-phase filter. Because the recording lacked an absolute timestamp, the audio was aligned empirically with the BCHH sequence. The approximately 366.5--366.8~m~s$^{-1}$ speed implicit in this display alignment was not interpreted as an atmospheric propagation measurement. Unless the earlier frame-by-frame analysis can be fully reconstructed---including camera locations, video identifiers, frame rates, synchronization procedures, and timing uncertainties---the video/audio alignment should be described as illustrative rather than as an independent absolute timing constraint.

\subsection{Polarization analysis and synchronized forensic visualization}
\label{sec:methods_polarization}

The original forensic analysis used a synchronized MATLAB/GISMO visualization
to examine the relationship between public video observations and the BCHH
seismo-acoustic waveforms. The display comprised four coordinated elements:
(1) time-referenced video of the accident; (2) scrolling waveforms from the
three seismic and three infrasound channels; (3) a magnified approximately
5-s view of the six-channel waveform data; and (4) time-dependent seismic
polarization attributes derived from the three-component broadband
seismometer. This display was used to examine the onset and principal
explosion, search for a possible payload- or capsule-impact signal, and
interpret the prolonged continuous signal following the initial failure phase.

The legacy polarization analysis included rectilinearity, planarity,
ellipticity, particle-motion orientation, and related measures derived from
the three-component covariance matrix within moving time windows. These
attributes were used to distinguish intervals dominated by approximately
linear body-wave motion from those exhibiting elliptical or planar motion
consistent with surface waves. The archived MATLAB results also formed the
basis for earlier interpretations of a possible impact-related seismic phase
and mixed wave types within the prolonged continuous signal.

Because the exact window lengths, overlap, filtering, eigenvalue conventions,
coordinate orientation, and attribute definitions used in the original
MATLAB implementation have not yet been fully recovered, those interpretations
are treated as provisional in the present study. The polarization analysis
will be reimplemented in Python using the response-corrected three-component
ground-velocity records. The Python implementation will document all processing
parameters explicitly and will be treated as authoritative for the final
polarization results and figures.

The regenerated products will include the time-dependent polarization
attributes, particle-motion plots for selected intervals, and a synchronized
four-panel video combining the public accident video, scrolling six-channel
waveforms, a magnified waveform window, and the polarization time series.
The synchronized visualization will be distributed as supplementary material.

\subsection{Interpretation and provenance}

Throughout the analysis, direct observations were distinguished from later physical interpretations. Direct observations included arrival times, waveform amplitudes, pressure polarity and duration, apparent propagation speeds, back azimuths, coherence measures, ground velocities, and derived energy estimates. Interpretations of individual signals were based on their temporal correspondence with publicly available video and, where relevant, information released during the subsequent SpaceX failure investigation.

The modern Python workflow was treated as authoritative for response
correction, baseline removal, event association, pulse measurements, planar-
and circular-wavefront analysis, acoustic and seismic amplitude measurements,
and amplitude regression. Python results take priority over legacy MATLAB
values wherever both are available. The polarization and particle-motion
analysis is still being reimplemented in Python; until it has been validated,
legacy polarization-derived interpretations of the payload- or capsule-related
interval and the prolonged continuous signal remain provisional. Legacy MATLAB
results were otherwise retained only where the corresponding analysis has not
yet been reproduced in Python, particularly for acoustic and seismic energies
and equivalent magnitudes. Where both workflows produced related quantities,
their definitions and provenance were kept explicit rather than silently
combining the results.

\subsection{Reproducibility}

Each Python processing stage writes its numerical products, configuration parameters, intermediate event streams, event-selection decisions, and diagnostic figures to a dedicated output directory. The canonical workflow comprises notebooks for response correction, input preparation, meteorological processing, candidate association, planar-wave analysis, circular-wavefront analysis, acoustic--seismic amplitude analysis, and production of paper figures and tables. Event-filter outcomes are retained for all candidates, including events that fail the permissive or core thresholds.

The archive will also include the Python implementation of the
three-component polarization analysis, the resulting regularly sampled
polarization attributes with complete processing metadata, selected
particle-motion products, and the script used to generate the synchronized
four-panel supplementary video. The source video identifier, synchronization
offset, frame-rate handling, waveform time base, filtering, moving-window
parameters, coordinate conventions, and definitions of all polarization
attributes will be documented explicitly. Where redistribution of the original
public video is not permitted, the archive will provide the synchronization
metadata and rendering scripts together with an exact citation and identifier
for the source video rather than redistributing the source footage.

The continuous waveforms, station metadata, manual arrival picks, event catalogue, Python notebooks, legacy MATLAB routines, derived tables, and figure-generation scripts will be archived in an open repository with a persistent identifier. The archive will identify the specific workflow and software implementation used to generate each principal result. The Python processing uses ObsPy \citep{beyreuther_obspy_2010,krischer_obspy_2015}
and the scientific Python ecosystem, whereas the original analysis used
Antelope \citep{brtt_antelope}, GISMO \citep{thompson_gismo_2017}, and custom
MATLAB routines. Final software and data citations will also document NumPy,
SciPy, pandas, Matplotlib, pyproj, the KSC meteorological observations, the
instrument responses and infraBSU calibration information, and the public
videos used in the chronology.

\section{Results}
\label{sec:results}

\subsection{Atmospheric conditions}
\label{sec:results_atmosphere}

Atmospheric conditions remained nearly constant throughout the 30-minute explosion sequence. Averaged over all weather-tower observations, the air temperature was 26.36~\textdegree C and the relative humidity was 83.87\%. Vector averaging of the wind measurements yielded a mean along-path wind component of $+2.65~\mathrm{m\,s^{-1}}$. These conditions correspond to a still-air sound speed of $348.33~\mathrm{m\,s^{-1}}$ and an effective propagation speed of $350.98~\mathrm{m\,s^{-1}}$ along the source-to-BCHH path after accounting for the component of the wind in the propagation direction. Little systematic temporal variation was observed, supporting use of a single atmospheric model for the acoustic analyses presented below.


\subsection{Candidate-event catalogue and temporal evolution}
\label{sec:results_catalogue}

Manual review yielded 384 accepted infrasound picks on HD1--HD3. Correction for
the predicted differential travel times reduced the normal approximately
\SIrange{80}{90}{\milli\second} traversal time across the array and allowed the
picks to be associated using a \SI{0.040}{\second} reduced-time window. The
association produced 154 candidate events, of which 63 contained picks on all
three infrasound channels and 91 contained picks on two channels. Thirteen
accepted picks remained unassociated.

The candidate catalogue extended from 13:07:15.927 to 13:33:54.396 UTC, a span
of approximately \SI{26.64}{\minute}. The activity was strongly episodic rather
than uniformly distributed through this interval (Figure~\ref{fig:overview}). Using the four
display intervals defined in Figure~\ref{fig:overview}, 38 candidates occurred during Phase~I,
3 during Phase~II, 87 during Phase~III, and 6 during Phase~IV; the remaining 20
candidates fell outside the marked intervals. The opening sequence contained the
largest-amplitude signals, whereas Phase~III contained the greatest number of
catalogued impulses. A period of sustained seismic tremor followed the opening
sequence between approximately 13:08:48 and 13:12:00 UTC, with no comparably
continuous signal on the infrasound channels.

The positive-compression/negative-rarefaction algorithm returned a valid pair
for 153 of the 154 candidate events. The sole exception was the principal
explosion, whose waveform was too complex to be represented by a single short
positive--negative cycle. On the initial descriptive classification, 46 events
were labelled robust impulsive events, 82 probable events, 25 weak coherent
candidates, and one---the principal explosion---a complex or noncanonical
waveform. These labels describe waveform quality rather than distinct physical
source types.

Of the 154 candidates, 148 passed the permissive criteria for planar array
analysis, including the explicitly retained principal explosion. Seventy-six
events passed the more conservative peak-to-peak pressure and SNR thresholds and
formed the core planar set. Inspection of the underlying amplitude, SNR, and
coherence distributions did not reveal sharp discontinuities at either set of
thresholds; the core set is therefore a conservative subset of a continuous
quality distribution rather than a separately resolved event population.



\begin{figure*}[t]
  \centering
  \includegraphics[width=0.98\textwidth]{fig03_bchh_30min_overview_zp_false.pdf}
  \caption{Overview of the Falcon~9 accident sequence recorded at BCHH. The
  upper three traces show acoustic pressure on HD1--HD3 and the lower three
  traces show three-component ground velocity. The four displayed phase
  intervals and the prolonged seismic-tremor interval are indicated. The
  displayed phase intervals contain 38, 3, 87, and 6 associated candidates,
  respectively; 20 additional candidates fall outside them. Time is UTC on
  1 September 2016.}
  \label{fig:overview}
\end{figure*}

\subsection{Absence of detectable precursors and later activity}
\label{sec:results_onset}

To test whether a detectable energetic precursor preceded the visible failure,
the complete continuous dataset spanning approximately 37~h before the first
accident-related signal was reviewed for anomalous impulsive seismic or
infrasonic activity. No such precursor attributable to the accident was
identified. The pre-event record contained normal background seismicity,
cultural noise, and environmental signals, but no coherent impulsive arrival
from the direction of SLC-40 that could be associated with an external
energetic event.

This result places an observational constraint on detectable seismo-acoustic
precursors at BCHH. It does not identify the initiating engineering mechanism
and cannot exclude aseismic, non-impulsive, weak, or otherwise undetectable
processes within the launch vehicle.

No additional explosion-related impulsive arrivals were identified during the
10.5~h following the final event in Phase~4.


\subsection{Opening sequence and key accident stages}
\label{sec:results_principal}

The opening sequence contained four geophysically distinct signals associated
with stages identified in the video chronology: the initial upper-stage failure,
the principal lower-stage explosion, a capsule impact, and a subsequent
capsule-related explosion (Table~\ref{tab:principal_events} and Figure~\ref{fig:opening}). The source-time estimates were
13:07:12.080, 13:07:15.750, 13:07:24.600, and 13:07:25.150 UTC, respectively.

The initial upper-stage failure produced a median-stack positive pressure of
\SI{23}{\pascal}, a negative pressure of \SI{-21}{\pascal}, and a peak-to-peak
pressure of \SI{44}{\pascal}. Its three-component vector PGV was
$7.85\times10^{-5}\ \si{\metre\per\second}$. The best planar solution gave a
back azimuth of \ang{201.5} and an apparent speed of
\SI{355.5}{\metre\per\second}.

Approximately \SI{3.67}{\second} later, the principal explosion produced the
largest acoustic and seismic amplitudes in the catalogue (Figure~\ref{fig:principal}). Using the
aligned median-stack estimator applied consistently to all events, the positive
pressure was \SI{1353}{\pascal}, the negative pressure was
\SI{-84}{\pascal}, and the peak-to-peak pressure was \SI{1437}{\pascal}. The
positive pressure was approximately 59 times that of the initial upper-stage
signal. Assuming geometric spreading proportional to $1/r$, the stack-positive
pressure corresponds to \SI{1919}{\pascal} at \SI{1}{\kilo\metre}. Its peak
sound-pressure level at BCHH was \SI{156.6}{\decibel} re
\SI{20}{\micro\pascal}.

Measurements made separately on the three infrasound sensors gave positive
peaks of 1.335, 1.546, and \SI{1.423}{\kilo\pascal} on HD1, HD2, and HD3,
respectively. The corresponding negative extrema were $-193$, $-231$, and
\SI{-234}{\pascal}. Thus, the median of the sensor-specific positive extrema was
\SI{1.423}{\kilo\pascal}, equivalent to approximately
\SI{2.014}{\kilo\pascal} at \SI{1}{\kilo\metre}. The difference from the
median-stack result reflects the use of two distinct estimators: the stack is
evaluated sample by sample after alignment, whereas the sensor-specific extrema
need not occur at exactly the same relative time.

The principal pressure waveform comprised an abrupt positive compression
followed by a longer, irregular decay and substantially smaller rarefaction. It
therefore did not display the compact, approximately symmetric positive--negative
morphology of many weaker catalogue events. Its shape qualitatively resembles a
blast pulse with an extended decay, although no parametric Friedlander-waveform
fit was performed. The seismic response persisted considerably longer than the
initial acoustic compression and contained large-amplitude motion on all three
components. The three-component vector PGV reached
$3.57\times10^{-3}\ \si{\metre\per\second}$, the largest value measured in the
catalogue.

The two capsule-related signals were separated by approximately
\SI{0.55}{\second}. The first produced stack pressures of $+276$ and
\SI{-102}{\pascal}, a peak-to-peak pressure of \SI{378}{\pascal}, and vector
PGV of $6.98\times10^{-4}\ \si{\metre\per\second}$. The second produced
$+291$ and \SI{-113}{\pascal}, a peak-to-peak pressure of
\SI{404}{\pascal}, and vector PGV of
$9.02\times10^{-4}\ \si{\metre\per\second}$. Their planar back azimuths were
\ang{201.5} and \ang{201.3}, and their apparent speeds were 351.5 and
\SI{351.0}{\metre\per\second}, respectively. Both directions and speeds were
consistent with propagation from the launch-complex sector.



\begin{figure*}[t]
  \centering
  \includegraphics[width=0.98\textwidth]{fig03_bchh_20s_zoom_with_audio_reduced_zp_false.pdf}
  \caption{Expanded view of the opening approximately 20~s of the accident
  sequence. Public-video audio is shown above the three infrasound and three
  seismic channels for illustrative synchronization. Markers identify the
  initial upper-stage failure, principal explosion, and two capsule-related
  signals after reduction for acoustic travel time. Because the video lacks an
  independent absolute timestamp, the audio alignment is illustrative.}
  \label{fig:opening}
\end{figure*}


\begin{figure*}[t]
  \centering
  \includegraphics[width=0.92\textwidth]{principal_explosion_waveforms.png}
  \caption{Calibrated pressure and three-component ground-velocity waveforms
  for the principal explosion. The aligned median-stack pressure estimator
  gives $+1353$~Pa, $-84$~Pa, and 1437~Pa peak-to-peak; the sensor-specific
  positive extrema are 1.335, 1.546, and 1.423~kPa on HD1--HD3.}
  \label{fig:principal}
\end{figure*}


\begin{figure*}[t]
  \centering
  \includegraphics[width=0.92\textwidth]{capsule_explosion_waveforms.png}
  \caption{Calibrated waveforms for the two capsule-related signals separated
  by approximately 0.55~s. Their aligned-stack peak-to-peak pressures are 378
  and 404~Pa, and their vector PGVs are $6.98\times10^{-4}$ and
  $9.02\times10^{-4}$~m~s$^{-1}$, respectively. The physical interpretations
  of the two signals remain tied to the video chronology.}
  \label{fig:capsule}
\end{figure*}


\begin{table*}[t]
\caption{Notebook-derived measurements for four key signals in the opening
sequence. Pressure values are from the aligned median-stack estimator; PGV is
the three-component vector peak ground velocity. Source-time interpretations
follow the current video chronology and should retain appropriate timing and
interpretive uncertainty.}
\label{tab:principal_events}
\centering
\small
\begin{tabular}{lcccccc}
\hline
Signal & Source time (UTC) & $p_+$ (Pa) & $p_-$ (Pa) & $p_{\mathrm{p2p}}$ (Pa) & Vector PGV (m s$^{-1}$) & Planar BAZ / speed \\
\hline
Initial upper-stage failure & 13:07:12.080 & 23 & $-21$ & 44 & $7.85\times10^{-5}$ & 201.5\textdegree{} / 355.5 m s$^{-1}$ \\
Principal explosion & 13:07:15.750 & 1353 & $-84$ & 1437 & $3.57\times10^{-3}$ & 200.2\textdegree{} / 352.5 m s$^{-1}$ \\
Capsule-related signal 1 & 13:07:24.600 & 276 & $-102$ & 378 & $6.98\times10^{-4}$ & 201.5\textdegree{} / 351.5 m s$^{-1}$ \\
Capsule-related signal 2 & 13:07:25.150 & 291 & $-113$ & 404 & $9.02\times10^{-4}$ & 201.3\textdegree{} / 351.0 m s$^{-1}$ \\
\hline
\end{tabular}
\end{table*}

\subsection{Seismo-acoustic waveform characteristics}
\label{sec:results_waveforms}

Many of the stronger impulsive events were recorded simultaneously by the
infrasound array and broadband seismometer, illustrating efficient coupling of
atmospheric pressure waves into the ground. The seismic response was dominated by ground-coupled airwaves followed by Rayleigh-wave energy that persisted for several seconds after passage of the direct acoustic wave. No convincing evidence was found for direct seismic body waves generated by most individual explosions, consistent with sources occurring primarily within the atmosphere.

The principal explosion produced the strongest seismic response, including high-amplitude Rayleigh waves that dominated ground motion for approximately 20~s. In contrast, many smaller impulsive events produced clear infrasound arrivals but only weak or undetectable seismic signals.

Public video shows the payload or capsule descending and apparently reaching
the ground approximately 11~s after the onset. No distinct direct seismic
impact arrival has yet been confirmed in the raw three-component waveforms,
which are dominated during this interval by ground-coupled airwaves. An
impact-related interpretation was proposed in the original MATLAB analysis
using rectilinearity, planarity, ellipticity, and particle-motion attributes,
but those calculations must be reproduced before the interpretation can be
evaluated further.

Following the initial phase of impulsive activity, a prolonged continuous
signal persisted for approximately 3.5~min. It is strongest on the seismic
channels but may also be present on HD2 and therefore should not presently be
classified as seismic-only. Preliminary comparison with video suggests a
possible association with sustained launch-pad fire, but this interpretation
requires verification using synchronized video and a consistent examination of
all calibrated channels.

The infrasound waveforms were predominantly impulsive, with typical rise times of approximately 0.01--0.02~s. Even for the largest explosion, however, the waveforms did not exhibit the classical symmetric N-wave morphology commonly associated with idealized blast waves.


\subsection{Polarization and particle-motion results}
\label{sec:results_polarization}

% TODO: populate after the legacy MATLAB polarization workflow has been
% reproduced and validated in Python.
Time-dependent rectilinearity, planarity, ellipticity, polarization direction,
and particle-motion attributes will be used to reassess two interpretations
from the original forensic analysis: whether a distinct direct seismic phase
is associated with the apparent payload or capsule ground contact, and which
wave types contribute to the prolonged continuous signal following Phase~I.
No quantitative polarization results are reported here until the analysis has
been reproduced using the documented Python workflow.


\subsection{Planar-wave array results}
\label{sec:results_array}

The 148 analyzed events produced a narrow distribution of best-fitting source
directions. The median planar back azimuth was \ang{200.4}, with an
interquartile range (IQR) of \ang{1.3}. The complete analyzed range was
\SIrange{195.8}{208.2}{\degree}, although the extremes were dominated by weak or broad
coherence surfaces. Within the 76-event core set, the median was \ang{200.5},
the IQR was \ang{1.2}, and the range was \SIrange{196.8}{205.3}{\degree}. The core
median lay approximately \ang{1.1} clockwise from the surveyed direction to
SLC-40 (\ang{199.4}). At a range of \SI{1.42}{\kilo\metre}, an angular
difference of \ang{1.1} corresponds to approximately \SI{27}{\metre} in the
cross-range direction, but this offset is not uniquely attributable to source
position because geometry, atmospheric structure, waveform differences, and
finite-range effects can produce comparable shifts.

The median apparent speed for all analyzed events was
\SI{352.0}{\metre\per\second}, with an IQR of
\SI{12.6}{\metre\per\second}. The core set was more tightly concentrated, with
a median of \SI{351.25}{\metre\per\second} and an IQR of
\SI{6.75}{\metre\per\second}. Both medians closely matched the
meteorologically estimated effective speed of
\SI{350.98}{\metre\per\second}. The wider complete-catalogue range of
\SIrange{280}{440}{\metre\per\second} included solutions that reached the
search boundaries and was associated mainly with weak or poorly constrained
events.

The median joint coherence score was 0.819 for all analyzed events and 0.942 for
the core set. Fifty planar solutions met the high-quality coherence and
window-stability criteria, 46 were classified as probable, and 52 were weak or
broad. Among the core events, 47 were high quality, 21 probable, and only 8 weak
or broad. Perturbing the scoring windows produced median standard deviations of
\ang{0.29} in back azimuth and \SI{1.89}{\metre\per\second} in speed across all
analyzed events. For the core set these decreased to \ang{0.24} and
\SI{1.44}{\metre\per\second}, respectively.

The pairwise results supported the visual observation that HD2 and HD3 were
usually more similar to one another than either was to HD1. Across all analyzed
events, the median HD2--HD3 correlation was 0.928, compared with 0.645 for
HD1--HD2 and 0.644 for HD1--HD3. Within the core set, the corresponding medians
were 0.973, 0.886, and 0.863. The median HD1 correlation deficit---the HD2--HD3
correlation minus the mean of the two HD1 pair correlations---was 0.197 for all
analyzed events and 0.068 for the core set. HD1 therefore reduced coherence most
strongly among the weaker events, while remaining well correlated for most core
events.

The principal explosion gave a planar back azimuth of \ang{200.2}, an apparent
speed of \SI{352.5}{\metre\per\second}, and a coherence score of 0.9989. Its
solution was exceptionally stable under changes to the scoring window. The
composite coherence surface for the core catalogue also formed a compact maximum
near the surveyed launch-complex direction and the meteorologically predicted
speed, demonstrating that the catalogue-scale result was not controlled solely
by the principal explosion.


\subsection{Finite-range circular-wavefront results}
\label{sec:results_finite_range}

Replacing the plane-wave approximation with exact horizontal distances to a
source fixed at SLC-40 produced only modest changes. The median preferred
circular-wavefront speed was \SI{355.75}{\metre\per\second}, compared with
\SI{352.0}{\metre\per\second} for the planar model. The circular model's median
coherence score was lower than the planar optimum by only 0.0047. This small
difference is expected because the planar model can vary both direction and
speed, whereas the fixed-source calculation constrains the source direction.

For the principal explosion, the fixed-SLC-40 circular solution preferred
\SI{355.0}{\metre\per\second} and gave a coherence score of 0.9965, compared
with \SI{352.5}{\metre\per\second} and 0.9989 for the planar solution. Thus,
wavefront curvature across the approximately \SI{30}{\metre} aperture did not
materially alter the inferred propagation characteristics.

The range-profile analysis demonstrated that the array could not determine
source range. At a coherence-score reduction of 0.01, the along-range profile
remained open at the search boundary for all 148 events, and 66 event maxima lay
on an along-range boundary. Cross-range behavior was better constrained but
remained broad: only 7 profiles were open at the cross-range boundaries, the
median best cross-range offset was \SI{25}{\metre}, and the median full width at
a score reduction of 0.01 was approximately \SI{115}{\metre}. The apparent
clustering of trial locations away from SLC-40 therefore does not constitute a
resolved source-location result. It primarily reflects a broad trade-off among
cross-range position, propagation speed, and small systematic delay errors.


\subsection{Event amplitudes through time}
\label{sec:results_temporal}

Across all 154 candidates, median stack peak-to-peak pressure was
\SI{29.5}{\pascal}, with an interquartile range of
\SIrange{16.2}{51.3}{\pascal}. Median vector PGV was
$4.26\times10^{-5}\ \si{\metre\per\second}$, with an interquartile range of
$2.27\times10^{-5}$ to
$1.05\times10^{-4}\ \si{\metre\per\second}$. Both acoustic and seismic
amplitudes were largest during the opening sequence, declined rapidly after the
principal explosion and capsule-related signals, and increased again during the
dense Phase~III sequence. Later events were generally weaker than the
opening explosions, although several isolated impulses remained measurable late
in the record.

The principal explosion was the largest event by both pressure and vector PGV.
The next two candidates immediately following it reached peak-to-peak pressures
of approximately 699 and \SI{632}{\pascal} and vector PGVs of approximately
$2.78\times10^{-3}\ \si{\metre\per\second}$. This concentration of the largest
signals within a few seconds shows that the opening failure sequence dominated
the measured acoustic and seismic amplitudes, even though most catalogued events
occurred later.

The median dimensional pressure-to-PGV ratio for the full catalogue was
$5.70\times10^{5}\ \si{\pascal\per\metre\second}$. The ratio showed much more
scatter among events with vector PGV below approximately
$10^{-4}\ \si{\metre\per\second}$. This behavior coincided with decreasing
seismic SNR and reduced capture of the event maximum, indicating that at least
part of the increased low-amplitude scatter resulted from the seismic noise
floor rather than a demonstrable change in acoustic--seismic coupling.


\subsection{Acoustic pressure versus ground velocity}
\label{sec:results_regression}

Twenty-seven events met the joint acoustic-SNR, seismic-SNR, and window-capture
criteria for regression . Within this subset, peak-to-peak pressure and
vector PGV were strongly correlated in logarithmic space. The fitted power-law
exponent was
\begin{equation}
  b = 0.868 \pm 0.048,
\end{equation}
where the quoted uncertainty is the ordinary least-squares standard error. The
95\% event-bootstrap interval was 0.734--0.971. The regression explained 93.0\%
of the variance in $\log_{10}$ pressure ($R^2=0.930$).

The fitted exponent was slightly below unity, implying a modestly sublinear
relationship over the observed amplitude range. Nevertheless, the proportional
model
\begin{equation}
  p_{\mathrm{p2p}} = K\,\mathrm{PGV}
\end{equation}
with
\begin{equation}
  K = \SI{0.447}{\mega\pascal\per\metre\second}
\end{equation}
provided a nearly comparable description. The free power-law and proportional
models had base-10 logarithmic RMSE values of 0.128 and 0.146, respectively,
equivalent to typical multiplicative residual factors of approximately 1.34 and
1.40. The small difference in predictive scatter indicates that a nearly
constant pressure-to-PGV ratio remains a useful first-order description, even
though the unconstrained regression favored a sublinear exponent.

The principal explosion lay within the high-quality regression subset. Its
pressure-to-PGV ratio was approximately
$4.02\times10^{5}\ \si{\pascal\per\metre\second}$, comparable to the fitted
proportionality constant. The largest events therefore extended the general
acoustic--seismic amplitude trend rather than forming a clearly separate
population. In contrast, many of the weakest candidates departed substantially
from the fitted trend and should not be used to infer changes in energy
partitioning without a more explicit treatment of detection limits and
uncertainties in both measurements.


\section{Discussion}
\label{sec:discussion}

\subsection{Independent forensic reconstruction}
\label{sec:discussion_reconstruction}

% TODO Figure: integrated geophysical/video chronology after the video timing is finalized.
The BCHH observations provide an independent forensic record of the Falcon~9
accident. The original analysis was organized around five practical questions:
whether a detectable energetic precursor occurred; how many discrete events
followed and when; how large the principal events were; whether the impulsive
signals originated consistently from SLC-40; and whether their propagation
speeds were compatible with atmospheric airwaves. The seismic and infrasound
records address these questions without reliance on proprietary launch-vehicle
telemetry.

No anomalous impulsive precursor attributable to the accident was identified
during the preceding 37~h. The sequence subsequently produced 154 associated candidate infrasound
 events over approximately 26.64~min; 153 yielded standard pulse measurements,
 with the principal
explosion dominating both pressure and ground-motion amplitudes. High-quality
array solutions consistently pointed toward SLC-40, and their apparent speeds
were broadly compatible with direct acoustic propagation under the observed
atmospheric conditions.

The combined seismic, infrasound, and publicly available video observations
permit a detailed reconstruction of the temporal evolution of the Falcon~9
explosion sequence. A final integrated-chronology figure will be inserted after the waveform and video timing products are frozen. Although the datasets record
different physical manifestations of the accident, they are mutually
consistent and together provide a more complete chronology than any individual
dataset alone.

The earliest geophysical observations include a weak seismic feature followed approximately 2.3~s later by the first impulsive infrasound event. The origin and significance of the earlier seismic feature require final verification against the waveform records. No anomalous seismic or infrasonic activity was detected during the preceding 37~h, indicating that the observed sequence began abruptly. The first visible flash in public video coincides closely with the initial infrasound event after accounting for acoustic travel times, providing an independent synchronization between geophysical and optical observations.

Approximately 4~s later, the largest explosion produced by far the strongest acoustic and seismic signals recorded by BCHH. Public video shows this event corresponding to rapid expansion of the principal fireball and catastrophic destruction of the launch vehicle. Several seconds later, acoustic and ground-coupled signals occurred during
the interval in which video shows the payload or capsule descending and
apparently reaching the ground. The current raw waveforms do not isolate an
unambiguous direct seismic impact phase, so the association remains tentative
pending regeneration of the legacy polarization analysis.

Following the initial failure phase, activity evolved through three
additional phases of progressively smaller impulsive events separated by
intervals of greatly reduced activity. Between the first and second phases, a
prolonged continuous signal persisted for approximately 3.5~min. The signal is
prominent seismically and may also be present on HD2, but it does not appear as
a sequence of discrete impulsive infrasound events. A possible association with
sustained launch-pad fire will be tested against synchronized video.

Taken together, the observations demonstrate that the Falcon~9 accident was not a single explosion but a prolonged sequence of discrete energetic events extending over more than 26~min. The event catalogue provides a quantitative framework for examining that evolution and comparing it with subsequent engineering investigations.

\subsection{Comparison with the official failure investigation}
\label{sec:discussion_investigation}

At the time the original analyses were performed in 2016--2017, the engineering cause of the accident had not yet been established. SpaceX subsequently released results of its failure investigation, providing an independent chronology against which the geophysical observations can be evaluated.

According to that investigation, the initiating failure involved rapid release
of helium from a composite overwrapped pressure vessel immersed within the
second-stage liquid-oxygen tank. The resulting overpressure led to structural
failure of the second stage, followed by rapid destruction of the launch
vehicle and ignition of the propellant inventory. Although the geophysical
observations alone cannot identify these engineering mechanisms, they
independently document the timing and relative magnitudes of the principal
energetic events. The absence of a detectable preceding energetic signal is
consistent with an initiating process internal to the vehicle, but the BCHH
observations alone cannot identify that process or exclude every possible
external mechanism.

The initial infrasound event represents the first identifiable atmospheric manifestation of the accident and is consistent in timing with the earliest visible evidence of vehicle failure. Approximately 4~s later, the principal explosion produced the largest acoustic and seismic signals of the sequence and coincided with rapid growth of the fireball in public video. The close correspondence between the independently derived geophysical chronology and subsequently published engineering timeline supports the interpretation that the event catalogue accurately captures the evolution of the accident.

Importantly, the event catalogue, propagation measurements, amplitude estimates, and timing relationships were derived from the seismic and infrasound recordings independently of the later engineering interpretation. The SpaceX findings therefore provide an external comparison rather than the basis for the geophysical measurements.

\subsection{Seismo-acoustic characteristics of a near-field rocket explosion}
\label{sec:discussion_characteristics}

The Falcon~9 dataset provides a rare opportunity to examine a catastrophic technological explosion at a source distance of only approximately 1.4~km. At this range, direct acoustic arrivals dominate, minimizing uncertainties associated with long-range atmospheric propagation that complicate regional infrasound observations.

The explosion radiated energy much more efficiently into the atmosphere than into the ground. The principal explosion generated peak pressures exceeding 1400~Pa at BCHH while producing comparatively modest seismic amplitudes, yielding an acoustic-to-seismic energy ratio characteristic of a predominantly atmospheric source. Smaller events frequently generated clear infrasound signals with little or no measurable seismic counterpart, demonstrating that pressure measurements provide substantially greater sensitivity to this source class than ground-motion observations alone.

Array measurements further demonstrate that the impulsive events were remarkably consistent in both back azimuth and apparent propagation speed despite spanning more than two orders of magnitude in pressure amplitude. Nearly all high-quality events originated from the direction of SLC-40, indicating that the numerous impulsive sources occurred within or immediately adjacent to the launch complex.

Classical symmetric N-wave morphology was not observed, even for the
principal explosion. The available data do not uniquely distinguish among
source complexity, near-ground propagation, waveform evolution over the
1.4-km path, and instrumental or local-site effects, so no single explanation
is preferred.

\subsection{Planar and finite-range propagation}
\label{sec:discussion_propagation}

The modern Python analysis extends the original planar-wave treatment by explicitly evaluating a finite-range circular wavefront. At a source distance of approximately 1.4~km and an array aperture of approximately 30~m, wavefront curvature is small but potentially measurable. Comparison of the planar and fixed-source circular models provides a direct test of whether curvature affects the inferred propagation speed or coherence.

The finite-range sensitivity calculations also clarify the resolution of the three-element array. Two independent relative delays are sufficient to constrain direction and apparent speed under a planar model, but they cannot independently determine source easting, source northing, and propagation speed. The resulting open along-range coherence ridges therefore represent a fundamental range-resolution limitation rather than a failure of the optimization procedure. Final discussion of the magnitude of finite-range corrections will be added after the Python results are frozen.

\subsection{Comparison with other technological and natural explosions}
\label{sec:discussion_comparison}

Large technological explosions have previously been investigated using
seismic and infrasound observations, including the 2005 Buncefield oil-depot
explosion \citep{ottemoller_2008,green_2008,ceranna_2009}, the 2014 Moerdijk
chemical-plant explosions \citep{smink_2019}, and the 2020 Beirut
ammonium-nitrate explosion
\citep{pilger_beirut_2021,kim_pasyanos_2022}. Such studies demonstrate the value of seismo-acoustic observations for estimating source location, explosive yield, and source characteristics. Most, however, rely primarily on regional observations, where atmospheric propagation strongly influences arrival times, amplitudes, and waveforms and detailed propagation modelling is required.

The Falcon~9 dataset differs in several respects. BCHH was located only approximately 1.4~km from the launch pad, so direct acoustic and seismic arrivals dominate and propagation uncertainties are comparatively small. Rather than emphasizing yield estimation alone, the proximity of the observations permits reconstruction of a complex sequence comprising more than 150 individual impulsive events. Public video and the subsequent failure investigation provide independent constraints against which the geophysical chronology can be evaluated. Together, these datasets make the accident an unusually well-constrained benchmark for near-field seismo-acoustic analysis of a technological explosion.

In several respects, the prolonged sequence resembles sustained explosive activity during volcanic eruptions, in which a major explosion is followed by numerous smaller impulses over minutes to tens of minutes. The event-catalogue approach therefore has applications beyond rocket launches, providing a common framework for analysing complex atmospheric explosion sequences and integrating seismic, infrasound, and optical observations to reconstruct source evolution.

\subsection{Implications for rocket-launch monitoring}
\label{sec:discussion_monitoring}

The accidental recording of the Falcon~9 failure demonstrates the forensic
value of continuous seismo-acoustic monitoring at launch facilities. A single
well-calibrated station constrained the onset, absence of detectable energetic
precursors, number and timing of subsequent impulsive events, relative
amplitudes, source direction, and propagation characteristics without access to
proprietary engineering telemetry.

Since 2016, launch cadence and the number of launch providers at Kennedy Space Center and Cape Canaveral have increased substantially. The methods developed here---including systematic event cataloguing, array-based propagation analysis, and integration with independent observational datasets---provide a foundation for monitoring both nominal launches and off-nominal events.

Future networks employing multiple seismo-acoustic arrays could substantially improve source-location capability while enabling more detailed characterization of explosion dynamics, sonic booms, and other transient atmospheric phenomena associated with rocket operations.

\subsection{Limitations}
\label{sec:discussion_limitations}

Several limitations should be considered. The observations were obtained from a single three-element seismo-acoustic array, preventing independent triangulation of individual sources. The array analysis therefore provides robust directional and apparent-speed constraints but limited source-range resolution.

Event identification was performed manually because of the modest data volume and complexity of the sequence. Although this produced an internally consistent catalogue, larger datasets would benefit from automated detection algorithms trained and validated against the manual reference catalogue.

The infrasound calibrations were empirical field estimates rather than
laboratory-traceable absolute calibrations, and their absolute uncertainty
requires final documentation. Public-video timing is also limited by uncertain
camera locations, frame timing, and absent absolute timestamps in some
recordings.

The forensic conclusions are constrained by the types of processes detectable
seismo-acoustically. The absence of a precursor means that no anomalous
energetic seismic or infrasonic signal was identified above the effective noise
and detection limits of BCHH. It does not exclude slow structural degradation,
internal fluid or mechanical processes, weak sources, or other mechanisms that
would not generate a detectable signal at the array.

Interpretations of payload or capsule impact and of the prolonged continuous
signal remain provisional. The former depends partly on legacy polarization
attributes not yet regenerated in Python, while the latter requires
reassessment across the seismic and infrasound channels and comparison with
synchronized video.

Finally, the acoustic and seismic energy estimates should be regarded as
first-order approximations because of uncertainties in atmospheric propagation,
near-surface properties, source radiation, and the exact implementation of the
legacy MATLAB calculations. These limitations do not preclude relative
comparisons within the sequence but must be considered when comparing absolute
energies with other events.

\section{Conclusions}
\label{sec:conclusions}

The 1 September 2016 Falcon~9 static-fire accident produced one of the most comprehensively documented near-field technological explosions recorded by a colocated seismic and infrasound array. Continuous observations acquired at BCHH, approximately 1.4~km from SLC-40, captured the complete 26-minute sequence with high signal-to-noise ratio, providing a rare opportunity to investigate the temporal evolution of a catastrophic launch-vehicle failure using geophysical observations.

Manual association identified a catalogue of 154 candidate infrasound events.
Of these, 153 yielded valid positive-compression/negative-rarefaction
measurements; the principal explosion was retained as a complex noncanonical
event. The records also contain a prolonged interval of continuous seismic
tremor that is not counted as an additional discrete event. The interpretation of these features
requires further polarization and video analysis. The catalogue is the
principal scientific product of this study and provides quantitative
measurements of event timing, amplitude, apparent propagation speed, back
azimuth, and, where measurable, acoustic and seismic energy.

The observations address the principal forensic questions that motivated the
original analysis. No anomalous energetic seismic or infrasonic precursor
attributable to the accident was identified during the preceding 37~h. This
result constrains detectable impulsive precursors but does not exclude
non-energetic or seismo-acoustically undetectable initiating processes. The accident generated 154 associated candidate infrasound events over
approximately 26.64~min, rather than a single explosion. Approximately 3.67~s
after the initial upper-stage signal, the principal explosion produced the largest
pressure and ground-motion amplitudes of the sequence. High-quality array
solutions consistently pointed toward SLC-40, and their apparent propagation
speeds were broadly compatible with direct atmospheric propagation under the
observed meteorological conditions.

Comparison with public video and the subsequently published SpaceX failure investigation shows that the independently derived geophysical chronology is consistent with the documented evolution of the accident. The observations therefore provide an independent forensic record without relying on proprietary engineering telemetry.

More broadly, the study demonstrates the value of continuous near-field
seismo-acoustic monitoring at major launch facilities. Even a single
well-calibrated array was sufficient to constrain the timing, evolution,
relative magnitudes, source direction, and propagation characteristics of a
complex launch-vehicle failure. Multiple arrays would add source-location
capability and improve characterization of nominal launches, anomalous events,
sonic booms, and explosion dynamics.

By making the waveform archive, event catalogue, and analysis software openly available, this study will establish a benchmark dataset for developing and evaluating future seismo-acoustic methods and for comparison with subsequent rocket-launch accidents and other complex atmospheric explosions.

\begin{acknowledgements}
% TODO: confirm acknowledgements, funding, KSC support, array deployment assistance, weather-data providers, SpaceX wording, and Astronaut Beach House access.
Acknowledgements to be completed.
\end{acknowledgements}

\section*{Data and code availability}

% TODO: replace placeholders with persistent repository citations and DOIs.
Continuous waveform data, station metadata, manual arrival picks, the event catalogue, Python notebooks, archived MATLAB routines, derived tables, and figure-generation scripts will be deposited in an open repository with a persistent identifier. The repository will identify which workflow generated each reported product. GitHub-hosted development code will be archived as a versioned release in a persistent repository such as Zenodo.

\section*{Competing interests}

The authors declare that they have no competing interests.

% Requires mybibfile_updated.bib in the manuscript directory.
\bibliography{mybibfile}

\end{document}
